\section{问题重述}

随着社区老龄化程度加深，单纯依赖机构养老或家庭照护难以同时满足服务可达性、支付能力和运营效率要求。嵌入式社区养老服务站将助餐、日间照料、上门护理、康复理疗、助浴和紧急救助等服务嵌入居民生活圈，是提高社区养老服务供给能力的重要方式。赛题围绕若干小区的老人结构、服务需求、建设成本、距离关系和满意度规则，要求建立模型确定未来服务需求、服务站建设方案、收费补贴政策以及关键参数变化下的方案稳定性\cite{problem2026}。

本题研究对象包括 $10$ 个小区、三类老人、六类养老服务项目和三种服务站建设规模。三类老人分别为自理、半失能和失能老人；六类服务包括助餐、日间照料、上门护理、康复理疗、助浴和紧急救助；服务站规模分为小型、中型和大型。各类数据由题目附件给出，包括小区人口与老人结构、老人健康状态转移概率、服务需求频次、基准价格、直接支出、消费上限、服务站建设与运营成本、服务能力、小区间距离矩阵以及满意度评分规则。

围绕上述背景，需要依次解决以下四个问题。

第一，需要预测未来五年内各小区三类老人的数量变化，并据此测算第 $5$ 年末各类养老服务的理论需求。在此基础上，还需考虑老人收入和服务消费上限对收费服务需求的约束，得到消费约束后的实际服务需求。紧急救助具有公益免费属性，不应按收费服务口径削减，但仍需保留其服务量以便后续容量与成本核算。

第二，需要在建设预算和服务半径约束下确定服务站的数量、位置和规模。候选站点来自给定小区，服务站建设需要满足总建设预算不超过 $120$ 万元，小区到服务站的服务距离不超过 $1000$ 米。除地理覆盖外，还要结合满意度、服务能力、容量利用和需求满足程度评价方案优劣，并在确定方案后对基准价格下的财务表现进行测算。

第三，需要在问题二确定的服务站布局基础上，进一步制定服务价格和政府补贴方案。定价不仅影响老人支付负担，也会影响价格满意度、小区分配和最终有效服务人次，因此不能只作静态财务核算。模型还需计算服务站收入、直接成本、政府补贴、题目口径利润率和年度净收益，并在利润率不超过 $8\%$ 的约束下评价三类老人的经济、地理和服务满足可及性。

第四，需要分析关键参数变化对模型结果的影响。题目要求设置基准情景及若干扰动情景，包括老人增长率变化、健康状态转移概率变化、固定管理成本提高和建设预算调整。每个情景均应重新形成需求、选址、定价、补贴和财务结果，并比较站点方案、覆盖水平、满意度、利润率、年度净收益和补贴压力等指标，从而判断模型的灵敏度和方案的稳健性。
